%------------------------------------------------------------------------------%
\documentclass[a4paper,12pt,fleqn]{article}

\usepackage{calculo_varias_variaveis-1}

\ShowAnswers

%------------------------------------------------------------------------------%
\begin{document}

\makeHeader{CFVVI}{Prova 2 -- Turma 8}{04/04/2026}

\bigskip
\textbf{Atenção}: a aplicação das regras do produto ou do quociente em casos
nos quais um dos termos é constante \textbf{será considerada incorreta},
por evidenciar compreensão inadequada das regras de derivação.

% 01
%------------------------------------------------------------------------------%
\exercise{25\%}
Calcule o limite
\(
  \lim_{(x,y)\to(1,0)} \frac{(x-1)^2 - y^2}{(x-1)^2 + y^2},
\)
ou mostre que ele não existe

\begin{answer}
  \textbf{Atenção: os caminhos precisam passar pelo ponto $(1, 0)$}
  \bigskip

  Considerando o caminho $x=1$ temos
  \[
    g(y)
    = \frac{(x-1)^2 - y^2}{(x-1)^2 + y^2} \evalat{x=1}{}
    = \frac{(1-1)^2 - y^2}{(1-1)^2 + y^2}
    = \frac{-y^2}{y^2}
    = -1 \qquad \forall y\neq 0
  \]
  \[
    L_1
    = \lim_{y\to 0} g(y)
    = \lim_{y\to 0} (-1)
    = -1
  \]

  Considerando o caminho $y=0$ temos
  \[
    h(x)
    = \frac{(x-1)^2 - y^2}{(x-1)^2 + y^2} \evalat{y=0}{}
    = \frac{(x-1)^2 - 0^2}{(x-1)^2 + 0^2}
    = \frac{(x-1)^2}{(x-1)^2}
    = 1 \qquad \forall x\neq 1
  \]
  \[
    L_2
    = \lim_{x\to 1} h(x)
    = \lim_{x\to 1} 1
    = 1
  \]
  O limite não existe, pois assume valores diferentes em cada caminho
\end{answer}

% 02
%------------------------------------------------------------------------------%
\exercise{25\%}
Calcule $\DM{f}{x}{y}$ para
\(
  f(x, y) = y^2\sqrt{x^2+4} + x\ln(y) + e^{xy},
\)
na ordem indicada pela notação

\begin{answer}
  Como
  \[
    \DM{f}{x}{y} = \D{}{x} \left(\D{f}{y}\right)
  \]
  Primeiro calculamos a derivada de $f$ por $y$
  \begin{align*}
    \D{f}{y}
    & = \D{}{y} \left( y^2\sqrt{x^2+4} + x\ln(y) + e^{xy} \right)
    \\[2mm]
    & = \D{}{y} \left( y^2\sqrt{x^2+4} \right)
      + \D{}{y} \left( x\ln(y) \right)
      + \D{}{y} \left( e^{xy} \right)
    \\[2mm]
    & = \sqrt{x^2+4} \, \D{y^2}{y}
      + x \D{\ln(y)}{y}
      + e^{xy} \D{xy}{y}
    \\[2mm]
    & = \sqrt{x^2+4}\, 2y
      + x \frac{1}{y}
      + e^{xy} x
    \\[2mm]
    & = 2y \sqrt{x^2+4} + \frac{x}{y} + x e^{xy}
  \end{align*}
  Calculamos agora a derivada mista
  \begin{align*}
    \DM{f}{x}{y}
    & = \D{}{x} \left(\D{f}{y}\right)
    \\[2mm]
    & = \D{}{x} \left( 2y \sqrt{x^2+4} + \frac{x}{y} + x e^{xy}\right)
    \\[2mm]
    & = \D{}{x} \left( 2y \sqrt{x^2+4} \right)
      + \D{}{x} \left( \frac{x}{y}\right)
      + \D{}{x} \left( x e^{xy}\right)
    \\[2mm]
    & = 2y \D{}{x} \left(x^2+4\right)^{\sfrac{1}{2}}
      + \frac{1}{y}\D{x}{x}
      + \D{x}{x} e^{xy}
      + x \D{e^{xy}}{x}
    \\[2mm]
    & = 2y \frac{1}{2} \left(x^2+4\right)^{-\sfrac{1}{2}}
           \D{}{x} \left(x^2+4\right)
      + \frac{1}{y}
      + e^{xy}
      + x e^{xy} \D{xy}{x}
    \\[2mm]
    & = y \frac{1}{\sqrt{x^2+4}} 2x
      + \frac{1}{y}
      + e^{xy}
      + x e^{xy} y
    \\[2mm]
    & = \frac{2xy}{\sqrt{x^2+4}} + \frac{1}{y} + (1+xy)e^{xy}
  \end{align*}
\end{answer}

% 03
%------------------------------------------------------------------------------%
\exercise{25\%}
Use a regra da cadeia para calcular $\frac{dz}{dt}$, onde
\(
  z = \sqrt{1 + x^2 y\,},
\)
\(
  x = \cos\left(t^2\right),
\)
\(
  y = e^{2t}
\)

\begin{answer}
\begin{multicols}{2}
  Notamos que \(z = z(x, y)\), \(x = x(t)\) e \(y = y(t)\),
  portanto pela regra da cadeia temos
  \[
    \frac{dz}{dt}
    = \D{z}{x} \frac{dx}{dt}
    + \D{z}{y} \frac{dy}{dt}
  \]
  Calculando cada derivada
  \begin{align*}
    \D{z}{x}
    & = \D{}{x} \left( \sqrt{1 + x^2 y\,} \right)
    \\[2mm]
    & = \D{}{x} \left( 1 + x^2 y \right)^{\sfrac{1}{2}}
    \\[2mm]
    & = \frac{1}{2}\left( 1 + x^2 y \right)^{-\sfrac{1}{2}}
        \D{}{x} \left( 1 + x^2 y \right)
    \\[2mm]
    & = \frac{1}{2}\frac{1}{\sqrt{1 + x^2 y\,}} 2xy
    \\[2mm]
    & = \frac{xy}{\sqrt{1 + x^2 y}}
  \end{align*}
  \begin{align*}
    \D{z}{y}
    & = \D{}{y} \left( \sqrt{1 + x^2 y\,} \right)
    \\[2mm]
    & = \D{}{y} \left( 1 + x^2 y \right)^{\sfrac{1}{2}}
    \\*[2mm]
    & = \frac{1}{2}\left( 1 + x^2 y \right)^{-\sfrac{1}{2}}
      \D{}{y} \left( 1 + x^2 y \right)
    \\*[2mm]
    & = \frac{1}{2}\frac{1}{\sqrt{1 + x^2 y\,}} x^2
    \\*[2mm]
    & = \frac{x^2}{2\sqrt{1 + x^2 y}}
  \end{align*}

  \begin{align*}
    \frac{dx}{dt}
    & = \frac{d}{dt}\cos\left(t^2\right) \\[2mm]
    & = -\sin\left(t^2\right) \frac{dt^2}{dt} \\[2mm]
    & = -2t\sin\left(t^2\right)
  \end{align*}

  \[
    \frac{dy}{dt}
    = \frac{de^{2t}}{dt}
    = e^{2t}\frac{d}{dt}(2t)
    = 2e^{2t}
  \]

  \begin{align*}
    \frac{dz}{dt}
    & = \D{z}{x} \frac{dx}{dt}
      + \D{z}{y} \frac{dy}{dt}
    \\[2mm]
    & = \frac{xy}{\sqrt{1 + x^2 y}} \left(-2t\sin\left(t^2\right)\right)
      + \frac{x^2}{2\sqrt{1 + x^2 y}} \, 2e^{2t}
    \\[2mm]
    & = \frac{-2txy\sin\left(t^2\right)}{\sqrt{1 + x^2 y}}
      + \frac{x^2e^{2t}}{\sqrt{1 + x^2 y}}
    \\[2mm]
    & = \frac{x^2e^{2t}-2txy\sin\left(t^2\right)}{\sqrt{1 + x^2 y}}
    \\[2mm]
    & = \frac{\cos^2\left(t^2\right)e^{2t}-2t\cos\left(t^2\right)e^{2t}\sin\left(t^2\right)}
             {\sqrt{1 + \cos^2\left(t^2\right) e^{2t}}}
   \\[2mm]
   & = \cos\left(t^2\right)e^{2t}
     \frac{\cos\left(t^2\right) - 2t\sin\left(t^2\right)}
          {\sqrt{1 + \cos^2\left(t^2\right) e^{2t}}}
  \end{align*}
\end{multicols}
\end{answer}

% 04
%------------------------------------------------------------------------------%
\exercise{25\%}
Use derivação implícita para avaliar \(\frac{dy}{dx},\)
sabendo que a equação
\mbox{\(
  \sin(xy) + x^2 y = \cos(x+y),
\)}
define $y$ como função de $x$

\begin{answer}
  Vamos usar a fórmula
  \[
    \frac{dy}{dx} = - \frac{F_x}{F_y}
  \]
  com
  \[
    F(x, y) = \sin(xy) + x^2 y - \cos(x+y)
  \]

  Calculando as derivadas parciais
  \begin{align*}
    F_x = \D{F}{x} \;
    & = \D{}{x} \left(\sin(xy) + x^2 y - \cos(x+y)\right)
    \\[2mm]
    & = \D{}{x} \left(\sin(xy) \right)
      + \D{}{x} \left(x^2 y    \right)
      - \D{}{x} \left(\cos(x+y)\right)
    \\[2mm]
    & = \cos(xy)  \D{xy}{x}
      + y         \D{x^2}{x}
      + \sin(x+y) \D{}{x}(x+y)
    \\[2mm]
    & = y\cos(xy) + 2yx + \sin(x+y)
  \end{align*}

  \begin{align*}
    F_y = \D{F}{y} \;
    & = \D{}{y} \left(\sin(xy) + x^2 y - \cos(x+y)\right)
    \\[2mm]
    & = \D{}{y} \left(\sin(xy) \right)
      + \D{}{y} \left(x^2 y    \right)
      - \D{}{y} \left(\cos(x+y)\right)
    \\[2mm]
    & = \cos(xy)  \D{xy}{y}
      + x^2       \D{y}{y}
      + \sin(x+y) \D{}{y}(x+y)
    \\[2mm]
    & = x\cos(xy) + x^2 + \sin(x+y)
  \end{align*}
  Substituindo na fórmula temos
  \[
  \frac{dy}{dx}
  = - \frac{F_x}{F_y}
  = - \frac{y\cos(xy) + 2yx + \sin(x+y)}{x\cos(xy) + x^2 + \sin(x+y)}
  \]
\end{answer}

%------------------------------------------------------------------------------%
\end{document}
%------------------------------------------------------------------------------%
